\chapter{Introduzione all'Ingegneria del Software}

\section{Il corso in breve}

Prima di entrare nel vivo, conviene avere chiaro \textbf{cosa} si studierà e \textbf{come} si viene valutati,
perché il modo in cui è organizzato l'esame rispecchia molto da vicino i contenuti del corso.

Finora (Programmazione 1, Programmazione a Oggetti, \ldots) abbiamo imparato a scrivere \textit{programmi}.
In questo corso si fa un salto di scala: si impara a produrre \textit{software} in modo professionale,
cioè a lavorare su sistemi grandi, in gruppo, con clienti veri e con vincoli di tempo e di budget.
In inglese si parla di \hi{programming in the large}, in contrapposizione alla programmazione «in piccolo»
a cui siamo abituati.

\begin{table}[H]
\centering
\renewcommand{\arraystretch}{1.45}
\begin{tabularx}{0.92\textwidth}{>{\bfseries\leavevmode\color{primary!80!black}}P{3.6cm} L}
\toprule
Crediti e lezioni & 9 CFU, circa 36 lezioni tra teoria ed esercitazioni in aula. \\
Prerequisiti & Buona capacità di programmazione, progettazione e programmazione a oggetti, basi di UML (diagrammi delle classi e di sequenza). \\
Esame scritto & Domande ed esercizi numerici su tutto il programma. Può essere sostituito da \textbf{due prove intercorso} (una a metà corso e una a fine corso): serve la sufficienza in entrambe e il voto è la media. \\
Progetto di gruppo & Gruppi da 2 o 3 persone (più si è, più funzionalità vanno implementate). Tema unico con varianti per gruppo. \\
Voto finale & Scritto e progetto hanno \textbf{lo stesso peso}. La partecipazione attiva alle attività pratiche viene considerata positivamente. \\
Tempistiche & Scritto e discussione del progetto sono indipendenti e si possono sostenere in qualsiasi ordine. Chi consegna il progetto presto ha meno funzionalità da implementare. \\
\bottomrule
\end{tabularx}
\caption{Organizzazione del corso e modalità d'esame.}
\end{table}

\subsection{Il progetto come gioco di ruolo}

Il progetto non è un semplice «esercizio di programmazione», ma una \textbf{simulazione realistica} di come si lavora in azienda:

\begin{center}
\begin{tikzpicture}[node distance=1cm]
  \node[card=swblue, minimum width=5.2cm, minimum height=2.2cm] (cli) {\textbf{Cliente}\\[2pt]\small i docenti\\[2pt]\scriptsize(possono fornire requisiti\\\scriptsize incoerenti, come i clienti veri!)};
  \node[card=sworange, minimum width=5.2cm, minimum height=2.2cm, right=5.6cm of cli] (sh) {\textbf{Software house}\\[2pt]\small il gruppo di studenti\\[2pt]\scriptsize(responsabilità collettiva\\\scriptsize su tutto il progetto)};
  \draw[-{Stealth[length=3mm]}, line width=1.2pt, swblue] ([yshift=4mm]cli.east) -- node[above, font=\scriptsize, text=black]{requisiti, richieste, cambiamenti} ([yshift=4mm]sh.west);
  \draw[-{Stealth[length=3mm]}, line width=1.2pt, sworange] ([yshift=-4mm]sh.west) -- node[below, font=\scriptsize, text=black]{prodotto software + documentazione} ([yshift=-4mm]cli.east);
\end{tikzpicture}
\end{center}

Il gruppo deve replicare \textbf{l'intero processo software} (che vedremo nel Capitolo 3): ingegneria dei requisiti,
progettazione di sistema, progettazione degli oggetti e dell'interfaccia, implementazione e testing.
Quindi non si consegna (solo) codice, ma anche tutti i documenti che un'azienda produrrebbe.

La discussione del progetto avviene in tre momenti:
\begin{enumerate}
  \item \textbf{Presentazione tecnica} con slide (circa 10 minuti) per convincere i «clienti» della qualità del prodotto e delle scelte fatte;
  \item \textbf{Demo dal vivo} dell'applicazione su un portatile del gruppo;
  \item \textbf{Discussione} delle scelte di progettazione e implementazione.
\end{enumerate}

\warning{Plagio e contenuti generati}{
  Tutti gli elaborati consegnati vengono analizzati con strumenti automatici di rilevamento del plagio e dei contenuti generati
  da IA. Progetti troppo simili vengono invalidati e ai gruppi coinvolti viene assegnata una nuova traccia.
}

\subsection{Testi consigliati}

Le slide \textbf{non sostituiscono} lo studio sui libri. Per argomento, i testi suggeriti sono:

\begin{table}[H]
\centering
\renewcommand{\arraystretch}{1.3}
\begin{tabularx}{\textwidth}{>{\bfseries}P{3.8cm} L}
\toprule
Argomento & Testo \\
\midrule
Parti generali & I. Sommerville, \textit{Software Engineering}, Pearson; R. S. Pressman, \textit{Software Engineering: A Practitioner's Approach}, McGraw Hill \\
Progettazione a oggetti & C. Larman, \textit{Applicare UML e i Pattern}, Pearson \\
Codice & R. Martin, \textit{Clean Code}, Apogeo \\
UML & P. Stevens, R. Pooley, \textit{Usare UML}, Addison Wesley \\
Approfondimenti & P. Ammann, J. Offutt, \textit{Introduction to Software Testing}, Cambridge University Press; E. Gamma et al., \textit{Design Patterns}, Addison Wesley \\
Usabilità e HCI & B. Shneiderman et al., \textit{Designing the User Interface}, Pearson; J. Nielsen, \textit{Usability Engineering}, Morgan Kaufmann \\
\bottomrule
\end{tabularx}
\end{table}

Lettura consigliata: \textit{The Humble Programmer} di Edsger W. Dijkstra (1972), il discorso per il premio Turing
in cui Dijkstra riflette su come la crescita della potenza dei calcolatori abbia reso la programmazione
un problema \textit{più} difficile, non più facile.

% ============================================================
\section{Perché «ingegneria»?}

Partiamo da un'analogia. Osserviamo questi due ponti:

\begin{figure}[H]
\centering
\begin{minipage}[t]{0.47\textwidth}
  \centering
  \includegraphics[width=\linewidth, height=5.2cm, keepaspectratio]{ponte_artigianale}\\[4pt]
  \begin{tcolorbox}[colback=swlight, colframe=swgray, boxrule=0.5pt, arc=2pt, width=\linewidth]
  \small\textbf{Passerella in legno}\\
  1 artigiano, 10 giorni, circa 5.000\,€.\\
  Un approccio \textbf{non strutturato} è sufficiente.
  \end{tcolorbox}
\end{minipage}\hfill
\begin{minipage}[t]{0.47\textwidth}
  \centering
  \includegraphics[width=\linewidth, height=5.2cm, keepaspectratio]{ponte_san_giorgio_notte}\\[4pt]
  \begin{tcolorbox}[colback=swlight, colframe=swblue, boxrule=0.5pt, arc=2pt, width=\linewidth]
  \small\textbf{Viadotto San Giorgio (Genova)}\\
  1.000 persone, 330 aziende coinvolte.\\
  Serve un approccio \textbf{sistematico e strutturato}.
  \end{tcolorbox}
\end{minipage}
\caption{Stessa funzione (attraversare), complessità incomparabili.}
\end{figure}

Entrambi sono «ponti», ma nessuno costruirebbe un viadotto autostradale come si costruisce una passerella nel bosco:
servono progettisti, calcoli, normative, collaudi, coordinamento tra centinaia di persone.
Con il software succede esattamente la stessa cosa: un programmino da 200 righe si può scrivere «di getto»,
un sistema bancario da milioni di righe no.

\info{L'idea chiave}{
  L'\textbf{ingegneria} entra in gioco quando la complessità del problema supera la capacità di un singolo individuo di
  tenerlo sotto controllo. A quel punto l'improvvisazione non basta più: servono \textbf{metodi, processi e strumenti}.
}

% ============================================================
\section{Un po' di storia}

\begin{minipage}[t]{0.60\textwidth}
\vspace{0pt}
\subsection{Gli inizi: l'hardware è il protagonista}

Nel 1946 viene presentato l'\textbf{ENIAC} (\textit{Electronic Numerical Integrator and Computer}), uno dei primi calcolatori elettronici programmabili.
In quegli anni:
\begin{itemize}
  \item i calcolatori erano macchine enormi, costosissime e sbalorditive;
  \item la macchina era molto più «spettacolare» (e costosa) di qualche foglio di codice;
  \item la programmazione era considerata un'attività \textbf{secondaria} rispetto all'ingegneria dell'hardware;
  \item i programmi venivano sviluppati per lo più \textbf{individualmente}, spesso da chi poi li usava.
\end{itemize}
\end{minipage}\hfill
\begin{minipage}[t]{0.37\textwidth}
\vspace{0pt}
\centering
\includegraphics[width=\linewidth]{eniac}\\[2pt]
{\scriptsize Jean Bartik (a sinistra) e Frances Spence al lavoro sull'ENIAC.}
\end{minipage}

\begin{minipage}[t]{0.30\textwidth}
\vspace{0pt}
\centering
\includegraphics[width=0.95\linewidth]{margaret_hamilton}\\[2pt]
{\scriptsize Margaret Hamilton nel 1969, accanto al codice scritto dal suo team per l'Apollo Guidance Computer.}
\end{minipage}\hfill
\begin{minipage}[t]{0.66\textwidth}
\vspace{0pt}
\subsection{Anni '60: la crisi del software}

Negli anni '60 i sistemi diventano molto più complessi (si pensi ai programmi spaziali della NASA).
Gli approcci individuali \textbf{non scalano}: quello che funziona per un programma piccolo crolla su sistemi grandi.

Nasce così la cosiddetta \hi{crisi del software}: i sistemi erano
\begin{itemize}
  \item \textbf{inaffidabili} (pieni di errori);
  \item \textbf{fuori budget} (costavano molto più del previsto);
  \item \textbf{in ritardo} rispetto ai tempi di consegna pianificati.
\end{itemize}

Diventa evidente il bisogno di una nuova disciplina che studi \textbf{metodologie per la produzione del software}.
Il termine \textit{software engineering} si diffonde in quegli anni: fu usato da Margaret Hamilton,
che guidava il team del software di bordo delle missioni Apollo, e diede il titolo a una famosa conferenza NATO del 1968
dedicata proprio alla crisi del software.
\end{minipage}

% ============================================================
\section{L'approccio ingenuo}

L'approccio con cui abbiamo lavorato finora nei corsi di programmazione è più o meno questo:

\begin{center}
\begin{tikzpicture}
  \node[card=swteal, minimum width=3.2cm, minimum height=1.1cm] (req) {\textbf{Requisiti}\\\scriptsize(la traccia)};
  \node[single arrow, draw=swgray, fill=swgray!20, minimum height=3.4cm, minimum width=1.2cm, single arrow head extend=2mm, right=0.4cm of req, font=\small] (cod) {coding};
  \node[card=swteal, minimum width=3.2cm, minimum height=1.1cm, right=0.4cm of cod] (prog) {\textbf{Programma}};
  \node[font=\Large\bfseries, text=swred] at ([yshift=0.9cm]req.north) {?};
  \node[font=\Large\bfseries, text=swred] at ([yshift=0.9cm]cod.north) {?};
  \node[font=\Large\bfseries, text=swred] at ([yshift=0.9cm]prog.north) {?};
\end{tikzpicture}
\end{center}

Leggiamo la traccia, scriviamo il codice, otteniamo il programma. Funziona benissimo per un esercizio d'esame,
ma appena il problema diventa reale emergono subito delle domande a cui questo schema non sa rispondere:

\begin{table}[H]
\centering
\renewcommand{\arraystretch}{1.35}
\begin{tabularx}{\textwidth}{>{\bfseries}P{3.9cm} L}
\toprule
Aspetto & Domande senza risposta \\
\midrule
Requisiti & Da dove arrivano? Sono corretti? Sono completi? Il cliente sa davvero cosa vuole? \\
Organizzazione & Come strutturo il programma? In quali moduli e classi lo divido? \\
Lavoro in team & Come divido il lavoro tra più programmatori? Come ci coordiniamo? \\
Correttezza & I requisiti sono davvero soddisfatti? Come faccio a sapere che il programma è corretto? \\
Cambiamento & Cosa succede quando i requisiti cambiano a metà strada (e cambieranno)? \\
\bottomrule
\end{tabularx}
\end{table}

L'Ingegneria del Software esiste proprio per dare una risposta \textbf{sistematica} a ciascuna di queste domande.

% ============================================================
\section{La dimensione conta}

Qual è il programma più grande che abbiamo mai scritto? E prima ancora: \textbf{come si misura la «dimensione» di un software?}
Una misura semplice (anche se imperfetta) sono le \hi{LOC} (\textit{Lines Of Code}), cioè le righe di codice,
di solito contate escludendo commenti e righe vuote.

Un progetto medio del corso di Programmazione a Oggetti dello scorso anno contava circa \textbf{3.000 righe} di codice
non di commento (incluse quelle generate automaticamente) e circa \textbf{35 classi}. Sembra tanto lavoro, vero?
Confrontiamolo con alcuni sistemi reali:

\begin{figure}[H]
\centering
\begin{tikzpicture}
\begin{axis}[
  xbar, xmode=log, log origin=infty,
  width=0.66\textwidth, height=9.2cm,
  xmin=1000, xmax=2e10,
  bar width=9pt,
  y dir=reverse,
  ytick={1,...,9},
  yticklabels={{Progetto OOP (studenti)},{Space Shuttle},{Telescopio Hubble},{Rover Curiosity (Marte)},{Firefox},{Kernel Linux (2020)},{Windows XP},{Auto moderna di fascia alta},{Google (tutti i servizi)}},
  yticklabel style={font=\small},
  xlabel={Righe di codice (scala logaritmica)},
  xlabel style={font=\small},
  xtick={1e3,1e4,1e5,1e6,1e7,1e8,1e9,1e10},
  xticklabels={1K,10K,100K,1M,10M,100M,1G,10G},
  xticklabel style={font=\scriptsize},
  xmajorgrids, grid style={swgray!25},
  axis line style={swgray},
  nodes near coords, point meta=explicit symbolic,
  every node near coord/.append style={font=\scriptsize\bfseries, text=black!80},
  enlarge y limits=0.07,
]
\addplot[fill=swblue!70, draw=swblue] coordinates {
  (3000,1) [3 mila]
  (400000,2) [400 mila]
  (2000000,3) [2 milioni]
  (5000000,4) [5 milioni]
  (10000000,5) [10 milioni]
  (28000000,6) [28 milioni]
  (45000000,7) [45 milioni]
  (100000000,8) [100 milioni]
  (2000000000,9) [2 miliardi]
};
\end{axis}
\end{tikzpicture}
\caption{Ordini di grandezza del software reale (stime diffuse pubblicamente, arrotondate).}
\end{figure}

\warning{Attenzione alla scala}{
  L'asse orizzontale è \textbf{logaritmico}: ogni tacca vale 10 volte la precedente. Il progetto di OOP è quindi più piccolo
  di un fattore 100 rispetto al software dello Space Shuttle, e di quasi un milione di volte rispetto all'insieme dei servizi Google.
  A queste dimensioni nessuna persona può conoscere tutto il codice: \textbf{servono metodo e organizzazione}.
}

% ============================================================
\section{Programmi e prodotti software}

Una distinzione fondamentale è quella tra \textbf{programma} e \textbf{prodotto software}.

\dfn{Programma}{
  Software sviluppato generalmente da \textbf{una sola persona}, che spesso ne è anche l'utente principale.
  Non richiede un processo formale né documentazione e in genere non è commercializzabile.
}

\dfn{Prodotto software}{
  Software sviluppato da \textbf{team} e usato da \textbf{altre persone}. È un prodotto industriale, sviluppato e testato secondo
  standard, che comprende molto più del solo codice sorgente o dell'eseguibile e che deve \textbf{evolvere nel tempo}.
}

\begin{table}[H]
\centering
\renewcommand{\arraystretch}{1.35}
\begin{tabularx}{\textwidth}{>{\bfseries}P{3.3cm} L L}
\toprule
 & \textcolor{swteal}{Programma} & \textcolor{swblue}{Prodotto software} \\
\midrule
Chi lo sviluppa & Una persona & Uno o più team \\
Chi lo usa & Lo sviluppatore stesso & Altre persone (clienti, utenti finali) \\
Processo & Nessun processo formale & Processo strutturato e documentato \\
Costi & Bassi & Significativamente più alti \\
Cosa comprende & Codice (ed eseguibile) & Codice, documentazione, manuali utente, manuali di installazione e configurazione, suite di test automatici \\
Ciclo di vita & «Usa e getta», one-shot & Evolve nel tempo (versioni, aggiornamenti, correzioni) \\
Esempio & Uno script per rinominare le foto delle vacanze & Un'app di home banking \\
\bottomrule
\end{tabularx}
\caption{Programma e prodotto software a confronto.}
\end{table}

\subsection{Tipi di prodotti software}

I prodotti software si dividono in due grandi famiglie. La differenza chiave è \textbf{chi definisce i requisiti}.

\begin{figure}[H]
\centering
\begin{tikzpicture}
  % --- Riga 1: software generico
  \node[chip=swblue, minimum width=17.5cm, minimum height=0.75cm, font=\small\bfseries, anchor=west] at (0,0) {Software generico (\textit{off-the-shelf})};
  \node[card=swblue, minimum width=2.4cm, minimum height=1.1cm] (mk) at (1.4,-1.7) {Team di\\marketing};
  \node[card=swblue, minimum width=2.4cm, minimum height=1.1cm] (p1) at (5.2,-1.7) {Prodotto\\generico};
  \draw[flow] (mk) -- node[above, font=\scriptsize]{requisiti} (p1);
  \foreach \k/\dy in {1/0.55,2/0,3/-0.55}{
    \node[font=\scriptsize, anchor=west] (u\k) at (7.6,-1.7+\dy) {\iuser\ cliente \k};
    \draw[-{Stealth}, swgray] (p1.east) -- (u\k.west);
  }
  \node[note, anchor=west, text width=7.6cm] at (9.7,-1.7) {Pensato per un \textbf{mercato}: venduto o distribuito a chiunque sia interessato.\\[2pt]\textit{Esempi}: Microsoft Office, Photoshop, Firefox, un videogioco.};

  % --- Riga 2: software su commessa
  \node[chip=sworange, minimum width=17.5cm, minimum height=0.75cm, font=\small\bfseries, anchor=west] at (0,-3.4) {Software su commessa (\textit{customly-developed})};
  \node[card=sworange, minimum width=2.4cm, minimum height=1.1cm] (cl) at (1.4,-5.1) {\iuser\ Cliente};
  \node[card=sworange, minimum width=2.4cm, minimum height=1.1cm] (p2) at (5.2,-5.1) {Prodotto\\su misura};
  \draw[flow] ([yshift=2mm]cl.east) -- node[above, font=\scriptsize]{requisiti} ([yshift=2mm]p2.west);
  \draw[flow, sworange] ([yshift=-3mm]p2.west) -- node[below, font=\scriptsize]{consegna} ([yshift=-3mm]cl.east);
  \node[note, anchor=west, text width=7.6cm] at (9.7,-5.1) {Sviluppato per le \textbf{esigenze specifiche} di un committente, che definisce i requisiti.\\[2pt]\textit{Esempi}: il gestionale di un ospedale, il sistema paghe di un ente pubblico.};
\end{tikzpicture}
\caption{La differenza chiave: chi definisce i requisiti (il marketing o il cliente).}
\end{figure}

\begin{esercizio}{Programma o prodotto? Generico o su commessa?}
Classifica i seguenti casi.
\begin{enumerate}
  \item Uno script Python che uno studente usa per calcolare la media dei propri voti.
  \item Il sistema di prenotazione degli esami realizzato per un'università.
  \item Un videogioco venduto su una piattaforma di distribuzione digitale.
  \item Il software di gestione del magazzino commissionato da una catena di supermercati.
\end{enumerate}

\textbf{Soluzione.}
(1) \textit{Programma}: una persona, un utente, nessun processo.
(2) \textit{Prodotto su commessa}: i requisiti li detta l'università.
(3) \textit{Prodotto generico}: i requisiti li definisce l'azienda (marketing) pensando a un mercato.
(4) \textit{Prodotto su commessa}.
\end{esercizio}

% ============================================================
\section{Lo sviluppo professionale}

Lo sviluppo professionale del software non è (solo) un'\textbf{arte} artigianale né (solo) una \textbf{scienza}: è un \hi{processo industriale}.
\begin{itemize}
  \item Il software viene sviluppato da \textbf{aziende}, organizzate in \textbf{team} con molti membri.
  \item Si lavora con \textbf{vincoli di tempo e di costo} e con \textbf{requisiti di qualità} stringenti.
  \item Come per ogni processo industriale, sono state sviluppate \textbf{metodologie} che guidano progettazione, sviluppo e verifica.
        Queste metodologie non nascono dal nulla: \textbf{formalizzano l'esperienza} accumulata in decenni di progetti (riusciti e falliti).
\end{itemize}

\warning{Essere il miglior programmatore non basta}{
  Da professionisti dell'informatica queste metodologie vanno \textbf{conosciute}. Un ottimo programmatore che non sa lavorare
  dentro un processo, in un team, con un cliente, produrrà comunque software di scarsa qualità.
}

% ============================================================
\section{Cos'è l'Ingegneria del Software}

Esistono diverse definizioni, che guardano la disciplina da punti di vista leggermente diversi.

\dfn{Ingegneria del Software (Sommerville)}{
  È la \textbf{disciplina ingegneristica} che si occupa di \textbf{tutti gli aspetti} della produzione di software professionale.
}

Analizziamo le due parti in grassetto:
\begin{itemize}
  \item \textbf{Disciplina ingegneristica}: gli ingegneri \textit{fanno funzionare le cose}. Applicano teorie, metodi e strumenti
        \textbf{dove sono appropriati}, per risolvere problemi rispettando i vincoli organizzativi ed economici. Non si cerca la soluzione
        «perfetta» in astratto, ma quella migliore \textit{date le risorse disponibili}.
  \item \textbf{Tutti gli aspetti della produzione}: non solo la programmazione! Rientrano anche la \textbf{gestione del progetto}
        (pianificazione, persone, costi) e lo sviluppo di strumenti, metodi e teorie a supporto dello sviluppo.
\end{itemize}

\dfn{Ingegneria del Software (IEEE)}{
  È l'applicazione di un approccio \textbf{sistematico, disciplinato e quantificabile} allo sviluppo, all'esercizio e alla manutenzione
  del software; in altre parole, l'applicazione dell'ingegneria al software.
}

\dfn{Ingegneria del Software (Bruegge e Dutoit)}{
  È lo stato dell'arte dello sviluppo di \textbf{software di qualità}, \textbf{nei tempi} e \textbf{nei costi} previsti.
}

\info{Cosa hanno in comune le definizioni}{
  Tutte e tre sottolineano che non basta «far funzionare il codice»: conta \textbf{come} ci si arriva (in modo sistematico),
  \textbf{cosa} si ottiene (qualità) e \textbf{a quali condizioni} (tempo e budget). Inoltre il lavoro non finisce con la consegna:
  esercizio e manutenzione fanno parte della disciplina.
}

% ============================================================
\section{Perché è difficile}

Produrre software di alta qualità rispettando budget e vincoli organizzativi \textbf{non è facile}. I motivi principali sono tre.

\begin{center}
\begin{tikzpicture}
  \node[card=swblue, text width=5.3cm, minimum height=3.7cm, align=left] (a) {%
    \textbf{\large 1. È astratto e intangibile}\\[4pt]
    \small Non è vincolato dalle proprietà dei materiali o dalle leggi della fisica. Per questo può diventare
    \textbf{estremamente complesso} molto in fretta, senza che nessuno «veda» il problema.};
  \node[card=sworange, text width=5.3cm, minimum height=3.7cm, align=left, right=0.35cm of a] (b) {%
    \textbf{\large 2. È ovunque ed è diverso}\\[4pt]
    \small Infrastrutture nazionali, elettrodomestici, industria, finanza, intrattenimento\ldots\
    Ogni dominio ha esigenze molto diverse dagli altri.};
  \node[card=swred, text width=5.3cm, minimum height=3.7cm, align=left, right=0.35cm of b] (c) {%
    \textbf{\large 3. Non esiste un metodo universale}\\[4pt]
    \small Nessuna metodologia va bene per ogni tipo di software, azienda e situazione.
    E nessuna tecnologia è una \textbf{pallottola d'argento} che risolve tutto.};
\end{tikzpicture}
\end{center}

\info{No Silver Bullet}{
  L'espressione «pallottola d'argento» viene dal celebre articolo \textit{No Silver Bullet} di Fred Brooks (1986):
  come nelle leggende solo una pallottola d'argento uccide il lupo mannaro, così si sperava in una tecnologia capace di
  eliminare da sola la complessità del software. Brooks sosteneva che \textbf{non esiste}: la complessità è in buona parte
  \textit{essenziale}, cioè intrinseca al problema, e nessuno strumento può farla sparire.
}

\subsection{Progetti che sono andati male}

Ancora oggi i progetti software falliscono o incontrano problemi seri: sforano il budget, sforano i tempi di rilascio,
consegnano un prodotto che non rispetta i requisiti di qualità\ldots\ o che non funziona affatto.

\begin{table}[H]
\centering
\renewcommand{\arraystretch}{1.4}
\begin{tabularx}{\textwidth}{>{\bfseries}P{3.4cm} L >{\centering\arraybackslash}p{3.2cm}}
\toprule
Progetto & Cosa è successo & Costo \\
\midrule
Healthcare.gov (USA) & Portale sanitario commissionato dal governo federale. Nella prima settimana solo l'1\% degli utenti riuscì a iscriversi. & \textbf{1,5 mld \$}\newline\scriptsize(15 volte i 93,7 mln previsti) \\
Queensland Health (Australia) & Nuovo sistema paghe per 80.000 dipendenti di 16 dipartimenti. Avviato nel 2007 con rilascio previsto nel 2008: rilasciato nel 2010, funzionava a malapena, sistemato solo nel 2018. & \textbf{$\sim$850 mln \$}\newline\scriptsize(200 volte i 4,2 mln previsti) \\
Cyberpunk 2077 & Videogioco annunciato nel 2012, uscita prevista ad aprile 2020, rinviata più volte fino a dicembre 2020. Al rilascio era pieno di bug e ingiocabile su alcune console. & \textbf{$\sim$400 mln \$} \\
\midrule
\rowcolor{swlight} Per confronto: Viadotto San Giorgio & 1,1 km di lunghezza, 45 m di altezza, 4 corsie + 2 di emergenza. Costruito in circa \textbf{un anno}. & \textbf{202 mln €} \\
\bottomrule
\end{tabularx}
\caption{Alcuni progetti software problematici, confrontati con un'opera di ingegneria civile.}
\end{table}

Il confronto è impietoso: un viadotto autostradale è costato meno di molti sistemi informativi falliti.
L'ingegneria civile ha secoli di esperienza e metodi consolidati; l'ingegneria del software è molto più giovane.

\subsection{Quando i problemi non sono solo economici}

A volte le conseguenze di un software difettoso vanno ben oltre il budget:

\begin{itemize}
  \item \textbf{London Ambulance Service (1992)}: il nuovo sistema computerizzato di smistamento delle ambulanze andò in crisi poco dopo l'avvio,
        causando enormi ritardi nei soccorsi; ai disservizi furono attribuiti numerosi decessi.
  \item \textbf{Ariane 5, volo 501 (1996)}: il razzo europeo si autodistrusse circa 40 secondi dopo il lancio. Un modulo software riusato
        dall'Ariane 4 convertiva un numero in virgola mobile a 64 bit in un intero a 16 bit: con le traiettorie più veloci del nuovo razzo
        il valore non ci stava più (\textit{overflow}).
  \item \textbf{Boeing 737 MAX (2018--2019)}: due incidenti con 346 vittime complessive, legati al sistema automatico di controllo MCAS,
        che si basava sui dati di un singolo sensore.
\end{itemize}

\error{Il messaggio}{
  Il software controlla ospedali, aerei, centrali, banche. Un difetto di \textbf{processo} (requisiti sbagliati, test insufficienti,
  riuso senza verifica, cattiva comunicazione) può costare vite umane. Per questo servono metodi ingegneristici.
}

% ============================================================
\section{Il lavoro dell'ingegnere del software}

Riassumendo, un ingegnere del software deve:
\begin{itemize}
  \item adottare un approccio \textbf{sistematico e organizzato} alla produzione del software;
  \item usare \textbf{strumenti, tecniche e metodologie appropriati}, scelti in base al problema da risolvere,
        ai vincoli di sviluppo e organizzativi e alle risorse disponibili;
\end{itemize}
per produrre \textbf{software di alta qualità}\ldots
\begin{center}
\begin{tikzpicture}
  \coordinate (Q) at (0,3.1);
  \coordinate (B) at (-3.3,0);
  \coordinate (T) at (3.3,0);
  \fill[swblue!8] (Q) -- (B) -- (T) -- cycle;
  \draw[line width=1.4pt, swblue] (Q) -- (B) -- (T) -- cycle;
  \node[chip=swblue, font=\small\bfseries, minimum width=2.6cm] at (Q) {Alta qualità};
  \node[chip=sworange, font=\small\bfseries, minimum width=2.6cm] at (B) {\ldots entro il budget};
  \node[chip=swred, font=\small\bfseries, minimum width=2.6cm] at (T) {\ldots entro la scadenza};
  \node[font=\small\itshape, text=swgray, align=center] at (0,1.05) {\ldots mentre\\i requisiti cambiano!};
  \draw[decorate, decoration={snake, amplitude=1.2pt, segment length=7pt}, swgray, line width=0.8pt, -{Stealth}]
     (-4.8,1.6) -- node[above, font=\scriptsize, sloped]{cambiamenti} (-2.1,1.6);
  \draw[decorate, decoration={snake, amplitude=1.2pt, segment length=7pt}, swgray, line width=0.8pt, -{Stealth}]
     (4.8,1.6) -- node[above, font=\scriptsize, sloped]{cambiamenti} (2.1,1.6);
\end{tikzpicture}
\end{center}

\gbox{In sintesi}{azure-gradient-6!80!black}{
\begin{itemize}
  \item L'Ingegneria del Software nasce dalla \textbf{crisi del software} degli anni '60: gli approcci individuali non scalano.
  \item Un \textbf{prodotto software} non è un programma: è sviluppato da team, per altri, include documentazione e test, ed evolve.
  \item Il software è \textbf{difficile} da ingegnerizzare: intangibile, ovunque, senza metodi universali.
  \item L'obiettivo è produrre \textbf{qualità} rispettando \textbf{tempi} e \textbf{costi}, gestendo il \textbf{cambiamento}.
\end{itemize}
}
